%! AP Statistics — Unit 9, Lesson 9.2 — Warm-Up
\documentclass[10pt]{article}
\usepackage{apstats-article}
\usepackage{apstats-boxes}

\begin{document}

\pageheader{Unit 9, Lesson 9.2}{Warm-Up}
\namedateperiod

\vspace{0.15in}

A study of 15 randomly selected college students recorded the number of hours spent
studying per week ($x$) and the score on a standardized exam ($y$).
The LSRL is $\hat{y} = 42.3 + 3.8x$.

\begin{enumerate}

  \item \textbf{Interpret the slope} of the LSRL in context.

        \writelines{2}

  \item One student studied 10 hours per week and earned a score of 81.
        \textbf{Calculate the residual} for this student.

        \vspace{0.05in}
        $\hat{y} = 42.3 + 3.8($\blank{1.5cm}$) = $ \blank{2.5cm}

        \vspace{0.08in}
        Residual $= y - \hat{y} = $ \blank{2cm} $-$ \blank{2cm} $= $ \blank{2cm}

        \vspace{0.1in}

  \item In Lesson 9.1, you learned about the sampling distribution of the slope $b$.
        Explain what it would mean if the \textbf{true slope $\beta = 0$} for this
        population.

        \writelines{2}

\end{enumerate}

\end{document}
